% ====================================================================
%  International Journal of Speech Technology — Springer
% --------------------------------------------------------------------
%  Template note: IJST is part of the Springer journal portfolio and
%  uses the Springer Nature LaTeX template (sn-jnl). For submission,
%  replace \documentclass{article} with the journal-specific class:
%      \documentclass[lnumbered,referee]{sn-jnl}
%      \jyear{2026}
%  and use the \title{}, \author{}, \aff[]{}, \abs{}, \keywords{},
%  \section*{} macros from the Springer template. We use the plain
%  `article` class here so the draft compiles anywhere; the structure
%  below follows the journal's expected section order
%  (Abstract, Keywords, Introduction, Related Work, Methodology,
%  Experimental Results, Discussion, Conclusion, References).
%  Typical IJST paper length: 12-18 typeset pages.
% ====================================================================
%% STATUS 2026-07-03: Experimental design superseded by the project audit
%% (research_logs/2026-07-03-project-audit-sota-roadmap.md). This venue awaits
%% the P4 flagship cross-lingual (Sinhala<->Tamil) study. Benchmark v0 results
%% now exist (see research_logs/2026-07-03-*) and are quoted below as
%% preliminary context; final tables require the SLCeleb v1 benchmark and the
%% P4 runs (cross-lingual trials, DANN/Dual-LoRA arms).
% ====================================================================
\documentclass[11pt,a4paper]{article}

% Engine guard: fontspec under XeLaTeX/LuaLaTeX (as in the 2026-05-23
% builds); T1/Latin Modern fallback so the draft also compiles with
% pdfLaTeX on machines without xetex/luaotfload.
\usepackage{iftex}
\iftutex
  \usepackage{fontspec}
  \setmainfont{DejaVu Serif}
  \setsansfont{DejaVu Sans}
  \setmonofont{DejaVu Sans Mono}
\else
  \usepackage[T1]{fontenc}
  \usepackage{lmodern}
\fi

\usepackage[a4paper,margin=1in]{geometry}
\usepackage[hidelinks]{hyperref}
\usepackage{booktabs,tabularx,amsmath,amssymb,microtype,enumitem,parskip}
\usepackage{authblk}

\title{Cross-lingual speaker verification for Sinhala and Tamil:
a single-corpus benchmark with ablation isolation of transfer,
calibration, and language-aware regularisation}

\author[1]{[First Author]}
\author[1]{[Second Author]}
\affil[1]{[Department], [University], Sri Lanka. \\
\texttt{[email@institution]}}

\date{}

\begin{document}
\maketitle

\begin{abstract}
\noindent We report the first published speaker-verification benchmark
for Sri Lankan Sinhala and Tamil, languages absent from all public SV
corpora and under-evaluated by even multilingual self-supervised
encoders. Evaluation is two-tier: benchmark v0, a 527-speaker
read-speech collection drawn from OpenSLR SLR52 (478 Sinhala speakers)
and SLR65 (49 Tamil speakers) with 12{,}000 seed-reproducible trials
per language, and benchmark v1, built on SLCeleb, our 280-speaker
in-the-wild Sinhala/Tamil corpus (pending data placement). On this
benchmark we compare training configurations --- zero-shot transfer of
VoxCeleb-trained checkpoints, training-free backend adaptation
(adaptive symmetric score normalisation, PLDA, language-aware
calibration), and in-language fine-tuning --- and isolate each lever
in a single-feature ablation. Per-language (monolingual Sinhala, monolingual
Tamil) and cross-lingual (enrolment-test language mismatch) EERs are
reported alongside pooled metrics. We describe the corpus-preparation
pipeline, the deterministic reproducibility protocol, and the
open-source toolkit released to support reproducible Sri Lankan SV
research. The results characterise the cross-lingual EER gap for a
single trained model and quantify each lever's contribution to closing
it.
\end{abstract}

\noindent \textbf{Keywords:} speaker verification; cross-lingual
transfer; low-resource speech; Sinhala; Tamil; multi-task learning;
score normalisation; PLDA; ECAPA-TDNN.

\section{Introduction}
\label{sec:intro}

Speaker verification (SV) --- the task of deciding whether two
recordings come from the same speaker --- has achieved remarkable
accuracy on English benchmarks: state-of-the-art systems report
0.8--1.5\% EER on VoxCeleb1-O \cite{desplanques2020ecapa,
nagrani2017voxceleb}. The accuracy on lower-resource languages is
less well characterised, particularly for Sinhala and Tamil, the two
primary languages of Sri Lanka, which together account for
approximately [N=22] million speakers but appear in no publicly
available SV corpus.

The Sri Lankan deployment context exposes three additional sources of
difficulty beyond pure label scarcity: \emph{(i)} bilingual and
code-switched speech is normal in everyday usage, so the deployment
trial distribution includes pairs where enrolment and test languages
differ; \emph{(ii)} the dominant deployment channel is telephony, with
attendant codec and bandwidth-restriction artefacts; and \emph{(iii)}
labelled-speaker collection is expensive at the scale a single
research effort can muster, capping the corpus at the low hundreds of
speakers.

This paper makes three contributions:

\begin{enumerate}
    \item A reproducible two-tier Sinhala/Tamil SV benchmark ---
    v0: 527 speakers (478 Sinhala from OpenSLR SLR52, 49 Tamil from
    OpenSLR SLR65 read speech) with 12{,}000 trials per language;
    v1: SLCeleb, our 280-speaker in-the-wild corpus (pending) --- with
    monolingual, cross-lingual, and pooled trial lists, generated
    deterministically from a single random seed.
    \item Single-feature ablations isolating the contribution of
    cross-lingual fine-tuning, layer-wise LR decay, multi-task
    language-ID head, adaptive symmetric score normalisation, and
    PLDA scoring, with three-seed reproducibility.
    \item An open-source toolkit \cite{slspv2026repo} extending
    the upstream \texttt{voxceleb\_trainer} codebase with eleven new
    features and twenty-six bug fixes, each selectable via a single
    configuration key.
\end{enumerate}

The remainder of this paper is organised as follows.
Section~\ref{sec:related} situates the work in the broader SV and
cross-lingual speech-processing literature.
Section~\ref{sec:method} describes the corpus, training, and
evaluation protocol.
Section~\ref{sec:experiments} reports results.
Section~\ref{sec:discussion} discusses limitations and future work.

\section{Related work}
\label{sec:related}

\subsection{Speaker verification architectures and losses}

ECAPA-TDNN \cite{desplanques2020ecapa} is the de-facto reference
architecture for non-self-supervised SV, combining Res2Net-style
multi-scale 1D convolutions, squeeze-and-excite channel re-weighting,
multi-layer feature aggregation, and channel-attentive statistical
pooling. AAM-Softmax \cite{deng2019arcface} is the dominant
embedding loss, producing discriminative representations on the
hypersphere through a margin-augmented cosine softmax.

Self-supervised front-ends --- WavLM \cite{chen2022wavlm}, wav2vec~2
\cite{baevski2020wav2vec2}, mHuBERT-147 \cite{pratap2024mms} ---
provide multilingual speech-aware representations that benefit
low-resource SV but at significant parameter cost (95M--300M).
Our toolkit ships infrastructure for both ECAPA-TDNN (paper's primary
encoder) and SSL front-ends (future work).

\subsection{Cross-lingual and low-resource SV}

When the target-language SV corpus is small ($\leq$~10k utterances),
training from scratch overfits. The standard remedy --- cross-lingual
fine-tuning from a high-resource (typically English VoxCeleb1)
checkpoint with selective freezing and layer-wise learning-rate
decay --- is well established for ASR and translation but
under-documented for SV-specific fine-tuning, motivating our explicit
ablation.

\subsection{Score normalisation and generative backends}

Adaptive symmetric score normalisation (AS-Norm)
\cite{matejka2017analysis} and two-covariance PLDA
\cite{sizov2014unifying,kenny2010} are eval-time backends widely used
in NIST SRE pipelines. AS-Norm normalises trial scores against
per-trial-side cohort statistics; PLDA replaces cosine scoring with a
likelihood ratio under same- / different-speaker hypotheses. The two
compose: PLDA scores can themselves be AS-Norm'd.

\subsection{Multi-task and adversarial language-aware training}

A multi-task language-ID head trained jointly with the speaker loss
regularises the embedding to retain language-discriminative
structure. Adversarial counterparts (DANN \cite{ganin2015dann}) pursue
the opposite goal: language-invariant embeddings. Both appear in the
literature; we evaluate the multi-task variant as the primary lever
and document DANN as infrastructure-ready future work.

\section{Methodology}
\label{sec:method}

\subsection{Corpus and trial construction}
\label{sec:method:corpus}

Evaluation is two-tier. Benchmark v0 comprises 527 speakers: 478
Sinhala from OpenSLR SLR52 (crowdsourced read speech, 16~kHz) and 49
Tamil from OpenSLR SLR65 (studio read speech, resampled to 16~kHz at
load), keeping speakers with at least 10 utterances. Benchmark v1 is
built on SLCeleb, our 280-speaker in-the-wild Sinhala/Tamil corpus
(multi-genre, multi-session, with bilingual speakers), and is pending
data placement. Within each speaker we hold out 20\% of utterances for
trial-pair construction; the remainder constitutes the training pool
(31{,}049 training utterances at a 60-utterance-per-speaker cap for
the fine-tuning experiments).

v0 trial pairs are sampled at a 1:3 target-to-impostor ratio at three
operating points:
\begin{itemize}
    \item Monolingual Sinhala: 3{,}000 target + 9{,}000 impostor pairs;
    \item Monolingual Tamil: 3{,}000 + 9{,}000 pairs;
    \item Cross-lingual (enrolment in one language, test in the other,
    for speakers present in both languages): empty at v0 --- the
    OpenSLR corpora contain no bilingual speakers --- and populated at
    v1 by SLCeleb's bilingual speakers.
\end{itemize}

Pair sampling is reproducible from a single random seed
(\texttt{tools/sl\_dataprep.py --seed 42}).

\subsection{Encoder, loss, and training}
\label{sec:method:encoder}

The speaker encoder is ECAPA-TDNN \cite{desplanques2020ecapa} with
$C=1024$ channels, 80 log-mel filterbanks (with pre-emphasis 0.97),
embedding dimension 192, and channel-attentive statistical pooling.
The training loss is AAM-Softmax \cite{deng2019arcface}
($m=0.2$, $s=30$). Training augmentation reuses the MUSAN + RIR
simulation chain with a fixed augmentation-probability distribution
(0.30 clean / 0.20 reverb / 0.15 music / 0.15 speech / 0.20 noise).

\subsection{Cross-lingual fine-tune protocol}
\label{sec:method:finetune}

Configurations P1 and P2 (Section~\ref{sec:experiments}) initialise
from an ECAPA-TDNN checkpoint pretrained on VoxCeleb1
\cite{nagrani2017voxceleb}. During SL fine-tuning the mel front-end
(\texttt{torchfb}, \texttt{instancenorm}) is frozen, the global
learning rate is scaled by $0.1\times$ from the from-scratch value,
and layer-wise LR decay with $\gamma=0.9$ is applied across the
four ECAPA layers so that the layer closest to the speaker loss head
receives the full LR and the input-side layer receives
$0.9^3 \approx 0.73\times$ of it.

\subsection{Multi-task language-ID training}
\label{sec:method:langaux}

A linear classifier $W_{\mathrm{aux}}\in\mathbb{R}^{4\times 192}$
maps the speaker embedding (pre-reshape, one prediction per utterance)
to four language classes (Sinhala, Tamil, English, mix). The combined
training loss is
\begin{equation}
\mathcal{L}_{\mathrm{total}}
= \mathcal{L}_{\mathrm{speaker}}
+ \lambda \, \mathcal{L}_{\mathrm{lang}},
\quad \lambda = 0.3.
\end{equation}
Speakers absent from the lookup carry \texttt{ignore\_index}$=-1$ and
contribute no gradient. The lookup is generated automatically from
the corpus directory layout by \texttt{tools/sl\_dataprep.py}.

\subsection{Eval-time scoring backends}
\label{sec:method:scoring}

Three eval-time scoring backends are compared:
\begin{enumerate}
    \item Cosine similarity on L2-normalised embeddings.
    \item Two-covariance PLDA \cite{sizov2014unifying} with LDA
    pre-reduction to $d=200$ and length normalisation; closed-form
    fit on a held-out PLDA-train set extracted from the SL training
    data.
    \item AS-Norm \cite{matejka2017analysis} with top-$K=300$ cohort
    statistics on a 500-speaker in-domain cohort, layered on top of
    either backend (1) or (2).
\end{enumerate}

\subsection{Per-language and pooled evaluation}
\label{sec:method:eval}

EER and MinDCF (with $P_\mathrm{target}=0.05$, $C_\mathrm{miss}=1$,
$C_\mathrm{fa}=1$) are reported per trial list (Si, Ta, Cs) and on
the concatenation. The trainer evaluates all four operating points in
a single invocation. We report two EER definitions: the conservative
$\max(\mathrm{FPR},\mathrm{FNR})$ form historically used by the
upstream toolkit and the standard literature
$(\mathrm{FPR}+\mathrm{FNR})/2$ form. We use the latter as the
primary metric to align with NIST SRE convention.

\subsection{Reproducibility protocol}
\label{sec:method:repro}

All experiments use the deterministic-mode switch
(cuDNN deterministic, TF32 disabled,
\texttt{torch.use\_deterministic\_algorithms} enabled). Primary
configurations are run with three random seeds (42, 123, 7); we
report $\mathrm{mean} \pm \mathrm{std}$. Single-feature ablations
use seed 42 only. Per-feature contribution is measured as $\Delta$EER
when the named feature is disabled from the full stack.

\section{Experimental results}
\label{sec:experiments}

\subsection{Preliminary monolingual baselines (benchmark v0)}
\label{sec:exp:prelim}

%% Numbers below are from the 2026-07-03 benchmark v0 runs
%% (research_logs/2026-07-03-{zeroshot-baseline-table,p2-backend-adaptation,
%% p3-peft-finetuning}.md). They are monolingual context for this paper's
%% cross-lingual study; the P4 cross-lingual arms are pending.

\begin{table}[t]
\centering
\caption{Preliminary monolingual EER (\%) on benchmark v0 (527
read-speech speakers; 12{,}000 trials per language). Fine-tuned rows
are $\mathrm{mean} \pm \mathrm{std}$ over three seeds. Preliminary,
benchmark v0: read speech, near single-session, closed-set speakers
for the fine-tuned rows; v1 (SLCeleb) provides the open-set,
in-the-wild condition.}
\label{tab:prelim}
\begin{tabular}{lcc}
\toprule
\textbf{System} & \textbf{Si} & \textbf{Ta} \\
\midrule
Zero-shot SpeechBrain ECAPA-TDNN            & 4.78 & 3.21 \\
Zero-shot ReDimNet-B6                       & 2.71 & 1.48 \\
ReDimNet-B6 + AS-Norm (training-free)       & 2.07 & 0.90 \\
WavLM-base+ LoRA fine-tune (3 seeds)        & 3.73 $\pm$ 0.49 & 3.11 $\pm$ 0.61 \\
WavLM-base+ full fine-tune (3 seeds)        & \textbf{1.69 $\pm$ 0.06} & \textbf{1.41 $\pm$ 0.10} \\
\bottomrule
\end{tabular}
\end{table}

Table~\ref{tab:prelim} reports preliminary monolingual results on
benchmark v0. Two findings shape the cross-lingual study: \emph{(i)}
full fine-tuning beat LoRA at every seed on this closed-set benchmark,
contrary to the published PEFT-beats-full-FT prediction for language
shift --- the v1 open-set rerun arbitrates; and \emph{(ii)}
per-language calibration is mandatory: swapping a calibrator fitted on
one language onto the other inflates Cllr by up to $5.7\times$. The
cross-lingual (Cs) trial columns below remain pending: they require
the bilingual speakers that only benchmark v1 (SLCeleb) provides.

\subsection{Primary configurations}
\label{sec:exp:primary}

\begin{table}[t]
\centering
\caption{Equal-error rate (\%) on the SL-SPV evaluation set,
$\mathrm{mean} \pm \mathrm{std}$ over three random seeds.}
\label{tab:headline}
\begin{tabular}{lcccc}
\toprule
\textbf{Config} & \textbf{Pooled} & \textbf{Si} & \textbf{Ta} & \textbf{Cs} \\
\midrule
P0 from scratch  & [X.X $\pm$ Y.Y] & [X.X $\pm$ Y.Y] & [X.X $\pm$ Y.Y] & [X.X $\pm$ Y.Y] \\
P1 + fine-tune   & [X.X $\pm$ Y.Y] & [X.X $\pm$ Y.Y] & [X.X $\pm$ Y.Y] & [X.X $\pm$ Y.Y] \\
P2 full stack    & \textbf{[X.X $\pm$ Y.Y]} & \textbf{[X.X $\pm$ Y.Y]}
                 & \textbf{[X.X $\pm$ Y.Y]} & \textbf{[X.X $\pm$ Y.Y]} \\
\bottomrule
\end{tabular}
\end{table}

\subsection{Single-feature ablations}
\label{sec:exp:ablation}

\begin{table}[t]
\centering
\caption{Single-feature ablations from the full stack (P2), seed 42.
$\Delta$ = pooled-EER increase when the named feature is removed.}
\label{tab:ablation}
\begin{tabular}{lcc}
\toprule
\textbf{Ablation} & \textbf{Pooled EER \%} & $\boldsymbol{\Delta}$ \\
\midrule
P2 (all features)   & [X.X] & 0 \\
$-$ fine-tune       & [X.X] & $+$[Y.Y] \\
$-$ LLRD            & [X.X] & $+$[Y.Y] \\
$-$ lang-aux        & [X.X] & $+$[Y.Y] \\
$-$ AS-Norm         & [X.X] & $+$[Y.Y] \\
$-$ PLDA            & [X.X] & $+$[Y.Y] \\
\bottomrule
\end{tabular}
\end{table}

\subsection{Per-language characterisation}
\label{sec:exp:perlang}

[INSERT: per-language EER comparison figure — bar chart with three
groups (Si / Ta / Cs) and three bars per group (P0 / P1 / P2). The
expected pattern: Cs EER is the highest of the three; the
Si--Ta gap shrinks as features are added.]

\subsection{Cross-lingual EER gap}
\label{sec:exp:crosslingual}

[INSERT: discussion of the magnitude of (Cs $-$ pooled) EER, which is
the headline number for cross-lingual robustness. Smaller is better.
Compare across P0 / P1 / P2 to characterise which interventions
help close the cross-lingual gap.]

\section{Discussion}
\label{sec:discussion}

\subsection{Cross-lingual transfer is the dominant lever}

[NOTE: discuss after results land. Expected pattern: the largest
$\Delta$ in Table~\ref{tab:ablation} is the fine-tune ablation, by
a wide margin, indicating that cross-lingual transfer from English
contributes more than any in-domain regularisation technique.]

\subsection{AS-Norm and PLDA composition}

[NOTE: discuss whether the AS-Norm contribution remains positive
when PLDA is on, or whether the two redundantly model the same
residual variance.]

\subsection{Limitations}

Benchmark v0, while larger than initially planned (527 speakers), is
read speech from a single source per language and near single-session
per speaker, so absolute EERs are optimistic and channel diversity is
limited; its fine-tuning evaluation is also closed-set
(utterance-disjoint but speaker-overlapping). Benchmark v1 (SLCeleb,
280 in-the-wild, multi-session speakers) addresses both. Continued self-supervised
pretraining on unlabelled SL audio --- shown in the literature to add
5--15\% relative EER on top of off-the-shelf multilingual SSL ---
is documented in the toolkit as infrastructure-ready future work
(FEATURE-009).

\section{Conclusion}
\label{sec:conclusion}

We presented a reproducible two-tier speaker-verification benchmark
for Sri Lankan Sinhala and Tamil (v0: 527 read-speech speakers from
OpenSLR SLR52/SLR65; v1: SLCeleb, our 280-speaker in-the-wild corpus,
pending), the first such benchmark to our knowledge. Single-feature ablations isolate the contributions of
cross-lingual fine-tuning, layer-wise LR decay, multi-task
language-ID, AS-Norm, and PLDA. The open-source toolkit, corpus
preparation pipeline, and reproducibility checklist are released to
support downstream low-resource SV research, particularly on the
languages of the South Asian region for which no published baselines
currently exist.

\section*{Acknowledgements}
[Funding, computational resources, advisor acknowledgements.]

\section*{Declarations}
\noindent \textbf{Funding.} [statement] \\
\noindent \textbf{Competing interests.} The authors have no competing
interests to declare. \\
\noindent \textbf{Data availability.} The corpus-preparation pipeline
and trial lists are publicly released alongside this paper; the raw
corpus is available subject to [statement about data licensing /
ethics-board approval].

\bibliographystyle{plain}
\bibliography{../references}

\end{document}
